                             IMC 2015, Blagoevgrad, Bulgaria

                                              Day 2, July 30, 2015




Problem 6.     Prove that                          ∞
                                                   X            1
                                                         √             < 2.
                                                   n=1
                                                             n (n + 1)
                                                                        (Proposed by Ivan Krijan, University of Zagreb)
Solution.   We prove that
                                                     1        2   2
                                              √             <√ −√    .                                              (1)
                                                  n (n + 1)    n n+1
                 √
Multiplying by       n(n + 1), the inequality (1) is equivalent with
                                                            p
                                           1 < 2(n + 1) − 2 n(n + 1)
                                             p
                                           2 n(n + 1) < n + (n + 1)
which is true by the AM-GM inequality.
   Applying (1) to the terms in the left-hand side,
                                   ∞                ∞           
                                   X        1       X    2    2
                                       √          <     √ −√       = 2.
                                   n=1
                                         n (n + 1) n=1    n  n+1



Problem 7.     Compute                                           Z A
                                                         1               1
                                                     lim               A x dx .
                                                    A→+∞ A        1

                                                                        (Proposed by Jan ’ustek, University of Ostrava)
Solution 1.   We prove that
                                               1 A 1
                                                 Z
                                         lim         A x dx = 1 .
                                        A→+∞ A 1

   For A > 1 the integrand is greater than 1, so
                            1 A 1           1 A
                              Z               Z
                                                          1          1
                                  A x dx >        1dx = (A − 1) = 1 − .
                           A 1             A 1            A          A
    In order to nd a tight upper bound, x two real numbers, δ > 0 and K > 0, and split the interval
into three parts at the points 1 + δ and K log A. Notice that for suciently large A (i.e., for A > A0 (δ, K)
with some A0 (δ, K) > 1) we have 1 + δ < K log A < A.) For A > 1 the integrand is decreasing, so we can
estimate it by its value at the starting points of the intervals:
                                   Z A                   Z 1+δ        Z K log A       Z A        
                               1          1     1
                                         A dx =
                                          x                       +                +                  <
                               A    1           A            1          1+δ             K log A

                       1                               1                       1
                                                                                  
                        =  δ · A + (K log A − 1 − δ)A 1+δ + (A − K log A)A K log A <
                       A
                      1             1                 1
                                                                        δ            1
                  <       δA + KA log A + A · A
                                    1+δ             K log A   = δ + KA− 1+δ log A + e K .
                      A
   Hence, for A > A0 (δ, K) we have
                                       1 A 1
                                        Z
                                 1                                δ          1
                            1− <            A x dx < δ + KA− 1+δ log A + e K .
                                 A     A 1
Taking the limit A → ∞ we obtain
                                                  Z A                        Z A
                                           1                1          1               1            1
                               1 ≤ lim inf              A dx ≤ lim sup
                                                            x                        A x dx ≤ δ + e K .
                                    A→∞ A          1             A→∞ A         1

Now from δ → +0 and K → ∞ we get
                                          1 A 1                  1 A 1
                                             Z                    Z
                              1 ≤ lim inf       A dx ≤ lim sup
                                                  x                 A x dx ≤ 1,
                                    A→∞ A 1               A→∞ A 1

so lim inf A1 1A A x dx = lim sup A1 1A A x dx = 1 and therefore
             R     1                 R    1

      A→∞                      A→∞
                                                                Z A
                                                            1           1
                                                        lim           A x dx = 1 .
                                                       A→+∞ A     1

Solution 2.    We will1 employ l'Hospital's    rule. R
    Let f (A, x) = A x , g(A, x) = x1 A x , F (A) = 1A f (A, x)dx and G(A) = 1A g(A, x)dx. Since ∂A f and
                                        1                                   R                     ∂
 ∂
∂A
   g are continuous, the parametric integrals F (A) and G(A) are dierentiable with respect to A, and
                                        Z A                                  Z A
                  0                             ∂                1                 1 1 −1       1  1
                F (A) = f (A, A) +                f (A, x)dx = A A +                 A x dx = A A + G(A),
                                         1     ∂A                             1    x               A
and                                                    Z A       1    Z A
                           0                      ∂             AA        1 1 −1
                          G (A) = g(A, A) +         g(A, x)dx =     +      2
                                                                             A x dx =
                                              1 ∂A              A      1 x
                                                   A      1       1
                                    1      −1 1 −1        AA     AA         1
                                =A +A           A x     =     −         +       .
                                          log A       1    A    A log A log A
      Since limA→∞ A A = 1, we can see that limA→∞ G0 (A) = 0. Aplying l'Hospital's rule to lim G(A) we
                      1
                                                                                                 A          A→∞
get
                                               G(A)        G0 (A)
                                                  lim= lim        = 0,
                                           A→∞    A   A→∞     1
so                                                             
                                         0             1   G(A)
                                    lim F (A) = lim A A +         = 1 + 0 = 1.
                                   A→∞         A→∞          A
      Now applying l'Hospital's rule to lim F (A)
                                              A
                                                  we get
                                               A→∞
                                             Z A
                                      1                 1           F (A)       F 0 (A)
                                  lim              A x dx = lim           = lim         = 1.
                                 A→+∞ A       1                 A→∞   A    A→∞      1



Problem 8.     Consider all 2626 words of length 26 in the Latin alphabet. Dene the weight of a word as
1/(k + 1), where k is the number of letters not used in this word. Prove that the sum of the weights of all
words is 375 .
                                             (Proposed by Fedor Petrov, St. Petersburg State University)
Solution.  Let n = 26, then 375 = (n + 1)n−1 . We use the following well-known
   Lemma. If f (x) is a polynomial of degree at most n, then its (n + 1)-st nite dierence vanishes:
∆ f (x) := n+1          i n+1
                              f (x + i) ≡ 0.
  n+1
             P               
                i=0 (−1)   i
   Proof. If ∆ is the operator which maps f (x) to f (x + 1) − f (x), then ∆        is indeed (n + 1)-st power
                                                                              n+1

of ∆ and the claim follows from   the observation   that ∆  decreases the power  of a polynomial.
   In other words, f (x) = i=1 (−1)              f (x + i). Applying this for f (x) = (n − x)n , substituting
                             Pn+1       i+1 n+1
                                               
                                             i
x = −1 and denoting i = j + 1 we get
                                 n                            n  
                                                                      (−1)j
                                         
                           n
                                 X   n+1 j        n
                                                             X   n
                 (n + 1) =     (−1)        (n − j) = (n + 1)        ·       · (n − j)n .
                           j=0
                                      j+1                    j=0
                                                                 j    j+1
The j -th summand nj · (−1)         · (n − j)n may be interpreted as follows: choose j letters, consider all
                              j

                              j+1
(n − j)n words without those letters and sum up (−1)         over all those words. Now we change the order of
                                                           j

                                                        j+1
summation, counting at rst by words. For any xed word W with k absent letters we get kj=0 kj · (−1)
                                                                                            P             j

                                                                                                        j+1
                                                                                                             =
 1
    · kj=0 (−1)j · k+1      1
                               , since the alternating sum of binomial coecients kj=−1 (−1)j · k+1   vanishes.
     P                                                                            P                
k+1                j+1
                         = k+1                                                                  j+1
That is, after changing order of summation we get exactly initial sum, and it equals (n + 1) .  n−1




Problem 9.    An n × n complex matrix A is called t-normal if AAt = At A where At is the transpose of
A. For each n, determine the maximum dimension of a linear space of complex n × n matrices consisting
of t-normal matrices.
                                      (Proposed by Shachar Carmeli, Weizmann Institute of Science)
Solution.
   Answer  : The maximum dimension of such a space is n(n+1) 2
                                                                .
   The number 2 can be achieved, for example the symmetric matrices are obviously t-normal and
                  n(n+1)

they form a linear space with dimension n(n+1)
                                           2
                                               . We shall show that this is the maximal possible dimension.
   Let Mn denote the space of n × n complex matrices, let Sn ⊂ Mn be the subspace of all symmetric
matrices and let An ⊂ Mn be the subspace of all anti-symmetric matrices, i.e. matrices A for which
At = −A.
   Let V ⊂ Mn be a linear subspace consisting of t-normal matrices. We have to show that dim(V ) ≤
dim(Sn ). Let π : V → Sn denote the linear map π(A) = A + At . We have

                                   dim(V ) = dim(Ker (π)) + dim(Im (π))

so we have to prove that dim(Ker (π)) + dim(Im (π)) ≤ dim(Sn ). Notice that Ker (π) ⊆ An .
   We claim that for every A ∈ Ker (π) and B ∈ V , Aπ(B) = π(B)A. In other words, Ker (π) and Im (π)
commute. Indeed, if A, B ∈ V and A = −At then

                                (A + B)(A + B)t = (A + B)t (A + B) ⇔
                       ⇔ AAt + AB t + BAt + BB t = At A + At B + B t A + B t B ⇔
                       ⇔ AB t − BA = −AB + B t A ⇔ A(B + B t ) = (B + B t )A ⇔
                                         ⇔ Aπ(B) = π(B)A.

   Our bound on the dimension on V follows from the following lemma:
Lemma. Let X ⊆ Sn and Y ⊆ An be linear subspaces such that every element of X commutes with every
element of Y . Then
                                        dim(X) + dim(Y ) ≤ dim(Sn )

Proof.  Without loss of generality we may assume X = ZSn (Y ) := {x ∈ Sn : xy = yx ∀y ∈ Y }. Dene the
bilinear map B : Sn × An → C by B(x, y) = tr(d[x, y]) where [x, y] = xy − yx and d = diag(1, ..., n) is the
matrix with diagonal elements 1, ..., n and zeros o the diagonal. Clearly B(X, Y ) = {0}. Furthermore, if
y ∈ Y satises that B(x, y) = 0 for all x ∈ Sn then tr(d[x, y]) = −tr([d, x], y]) = 0 for every x ∈ Sn .
    We claim that {[d, x] : x ∈ Sn } = An . Let Eij denote the matrix with 1 in the entry (i, j) and 0 in
all other entries. Then a direct computation shows that [d, Eij ] = (j − i)Eij and therefore [d, Eij + Eji ] =
(j − i)(Eij − Eji ) and the collection {(j − i)(Eij − Eji )}1≤i<j≤n span An for i 6= j .
    It follows that if B(x, y) = 0 for all x ∈ Sn then tr(yz) = 0 for every z ∈ An . But then, taking z = ȳ ,
where ȳ is the entry-wise complex conjugate of y , we get 0 = tr(yȳ) = −tr(yȳt ) which is the sum of squares
of all the entries of y . This means that y = 0.
    It follows that if y1 , ..., yk ∈ Y are linearly independent then the equations
                                     B(x, y1 ) = 0,   ...,   B(x, yk ) = 0
are linearly independent as linear equations in x, otherwise there are a1 , ..., ak such that B(x, a1 y1 + ... +
ak yk ) = 0 for every x ∈ Sn , a contradiction to the observation above. Since the solution of k linearly
independent linear equations is of codimension k,
                                   dim({x ∈ Sn : [x, yi ] = 0, for i = 1, .., k}) ≤

                           ≤ dim(x ∈ Sn : B(x, yi ) = 0 for i = 1, ..., k) = dim(Sn ) − k.
The lemma follows by taking y1 , ..., yk to be a basis of Y .
   Since Ker (π) and Im (π) commute, by the lemma we deduce that
                                                                                         n(n + 1)
                         dim(V ) = dim(Ker (π)) + dim(Im (π)) ≤ dim(Sn ) =                        .
                                                                                            2



Problem 10.        Let n be a positive integer, and let p(x) be a polynomial of degree n with integer coecients.
Prove that
                                                                      1
                                                  max p(x) >             .
                                                 0≤x≤1                en
                                                             (Proposed by Géza Kós, Eötvös University, Budapest)
Solution.    Let
                                                 M = max p(x) .
                                                         0≤x≤1

   For every positive integer k, let                   Z 1
                                                                  2k
                                                Jk =          p(x) dx.
                                                         0
                                                                         2kn                       2kn
Obviously 0 < Jk < M 2k is a rational number. If (p(x))2k =                    ak,i xi then Jk =
                                                                                                         ak,i
                                                                                                                . Taking the least
                                                                         P                         P
                                                                                                         i+1
                                                                         i=0                       i=0
                                                          1
common denominator, we can see that Jk ≥                                .
                                             lcm(1, 2, . . . , 2kn + 1)
    An equivalent form of the prime number theorem is that log lcm(1, 2, . . . , N ) ∼ N if N → ∞. Therefore,
for every ε > 0 and suciently large k we have
                                       lcm(1, 2, . . . , 2kn + 1) < e(1+ε)(2kn+1)

and therefore
                                                          1                    1
                               M 2k > Jk ≥                              > (1+ε)(2kn+1) ,
                                             lcm(1, 2, . . . , 2kn + 1)  e
                                                                1
                                                M > (1+ε)(n+ 1 ) .
                                                        e          2k


Taking k → ∞ and then ε → +0 we get
                                                               1
                                                       M≥         .
                                                               en
Since e is transcendent, equality is impossible.


Remark. The constant
                          e ≈ 0.3679 is not sharp. It is known that the best constant is between 0.4213 and 0.4232.
                          1

(See I. E. Pritsker, The GelfondSchnirelman method in prime number theory, Canad. J. Math. 57 (2005),
10801101.)
